\subsection{The Day's Prayer}\label{sec:prayer}

In Advent the Church has her own words for every hour of the day, and a
Christian at home can join her at a few of them. What follows is a day in
Advent from morning to night. No one should try to take up all of it at once.
Choose one thing for the morning, one for the middle of the day and one for
the evening, and keep them.

\subsubsection*{Morning}

On the weekdays of Advent the Mass is that of the Sunday before, except on the
Ember Days, which have Masses of their own, and on a feast such as the
Immaculate Conception. So the Sunday's Collect belongs to the whole week. Say it
each morning. It is short enough to learn, and it puts the season's first word
into the first minutes of the day.

At Lauds, the morning Office, the Church sings a hymn that sounds like St
Paul's call to rise from sleep, \latin{En clara vox redarguit}, ``Lo, a clear voice
rebukes''. Its first lines in English run:

\begin{verse}
Hark! an awful voice is sounding; \\
``Christ is nigh!'' it seems to say; \\
``Cast away the dreams of darkness, \\
O ye children of the day!''
\end{verse}

\noindent The whole hymn is in Appendix~\ref{sec:prayers}, with the hymn of
Matins and the hymns of Christmas.

\subsubsection*{The Angelus}

The Angelus is the Gospel of the Ember Wednesday, the Annunciation, turned into
a prayer: the angel's message, the Virgin's consent, the Word made flesh, and a
collect that runs from the Annunciation to the Cross and the Resurrection. It is
said when the bell rings, morning, noon and evening. In Advent it means more
than at any other time, because the season is waiting for what it tells.

The day ends with the same words. At Compline, the Church's prayer before
sleep, she sings in Advent the antiphon \latin{Alma Redemptoris Mater}, ``Loving
Mother of the Redeemer'', and closes it with the versicle and prayer of the
Angelus. Both are in Appendix~\ref{sec:prayers}.

\subsubsection*{Reading}

Appendix~\ref{sec:daily} sets out for each day the chapter of Isaiah that the
Church reads at Matins, with a passage of the Fathers and a verse to keep
through the day. Read slowly. A short old treatise called the \latin{Scala
claustralium}, ``the ladder of the cloister'', describes four steps in such
reading: \latin{lectio, meditatio, oratio, et contemplatio}, reading,
meditation, prayer and contemplation. Read a few verses more than once; turn
them over until one phrase holds you; speak to God out of what it has shown
you; then be still before him. The treatise is ascribed to Guigo the
Carthusian.

St Alphonsus Liguori wrote eighteen meditations for Advent, and for the nine
days before Christmas, from 16 December, a novena of meditations. They may be
read beside the Fathers, or in their place.

\subsubsection*{Evening}

At Vespers the Church sings \latin{Creator alme siderum}, ``Kind Creator of the
stars'', in English ``Creator of the starry pole''. It moves from creation to
the Cross and from the Cross to the Judgment, and asks the Judge of the last
day for the weapons of grace. After the hymn, at the first Vespers of the
season, comes a versicle in Isaiah's words: ``Drop down dew, ye heavens, from
above, and let the clouds rain the just: let the earth be opened, and bud forth
a saviour''. From 17 December one of the great O antiphons is sung each evening
before and after the Magnificat; the seven are in Appendix~\ref{sec:prayers}.

Those who say the Little Office of Our Lady have its Advent form, which begins
with the season. Its Vespers antiphon is \latin{Missus est Gabriel}, ``Gabriel
was sent'', from the Gospel of the Annunciation.

There is also the chant \latin{Rorate caeli}, a prayer of the churches of
France in Advent: verses from Isaiah sung under the versicle's first words as a
refrain. ``Be not angry, O Lord, and remember no longer our iniquity''; ``We
have sinned, and we are become as one unclean''; ``Be comforted, be comforted,
my people; thy salvation shall speedily come''. It is in
Appendix~\ref{sec:prayers}, and it is good to sing or say at the evening meal in
the last week.

\subsubsection*{Night}

End the day with an account of it. St Ignatius of Loyola teaches the examen in
five points: give thanks to God for the benefits received; ask grace to know
your sins; go through the day, ``hour by hour, or period by period: and first as
to thoughts, and then as to words, and then as to acts''; ask pardon; resolve
to amend with his grace. In Advent examine yourself on the week's Gospel and
Epistle, with the questions of the examination by the weeks of Advent
(Section~\ref{sec:confession}).

Then the \latin{Alma}, or the O antiphon of the day, and sleep. The Church will
be watching through the night at Matins, and in the morning there is the
Collect again.
